% Template for PLoS
% Version 3.3 June 2016
%
% % % % % % % % % % % % % % % % % % % % % %
%
% -- IMPORTANT NOTE
%
% This template contains comments intended 
% to minimize problems and delays during our production 
% process. Please follow the template instructions
% whenever possible.
%
% % % % % % % % % % % % % % % % % % % % % % % 
%
% Once your paper is accepted for publication, 
% PLEASE REMOVE ALL TRACKED CHANGES in this file 
% and leave only the final text of your manuscript. 
% PLOS recommends the use of latexdiff to track changes during review, as this
% will help to maintain a clean tex file.  Visit
% https://www.ctan.org/pkg/latexdiff?lang=en for info or contact us at
% latex@plos.org.
%
%
% There are no restrictions on package use within the LaTeX files except that
% no packages listed in the template may be deleted.
%
% Please do not include colors or graphics in the text.
%
% The manuscript LaTeX source should be contained within a single file (do not
% use \input, \externaldocument, or similar commands).
%
% % % % % % % % % % % % % % % % % % % % % % %
%
% -- FIGURES AND TABLES
%
% Please include tables/figure captions directly after the paragraph where they
% are first cited in the text.
%
% DO NOT INCLUDE GRAPHICS IN YOUR MANUSCRIPT
% - Figures should be uploaded separately from your manuscript file. 
% - Figures generated using LaTeX should be extracted and removed from the PDF
%   before submission. 
% - Figures containing multiple panels/subfigures must be combined into one
%   image file before submission.
% For figure citations, please use "Fig" instead of "Figure".  See
% http://journals.plos.org/plosone/s/figures for PLOS figure guidelines.
%
% Tables should be cell-based and may not contain:
% - spacing/line breaks within cells to alter layout or alignment
% - do not nest tabular environments (no tabular environments within tabular
%   environments)
% - no graphics or colored text (cell background color/shading OK)
% See http://journals.plos.org/plosone/s/tables for table guidelines.
%
% For tables that exceed the width of the text column, use the adjustwidth
% environment as illustrated in the example table in text below.
%
% % % % % % % % % % % % % % % % % % % % % % % %
%
% -- EQUATIONS, MATH SYMBOLS, SUBSCRIPTS, AND SUPERSCRIPTS
%
% IMPORTANT Below are a few tips to help format your equations and other
% special characters according to our specifications. For more tips to help
% reduce the possibility of formatting errors during conversion, please see our
% LaTeX guidelines at http://journals.plos.org/plosone/s/latex
%
% For inline equations, please be sure to include all portions of an equation
% in the math environment.  For example, x$^2$ is incorrect; this should be
% formatted as $x^2$ (or $\mathrm{x}^2$ if the romanized font is desired).
%
% Do not include text that is not math in the math environment. For example,
% CO2 should be written as CO\textsubscript{2} instead of CO$_2$.
%
% Please add line breaks to long display equations when possible in order to
% fit size of the column. 
%
% For inline equations, please do not include punctuation (commas, etc) within
% the math environment unless this is part of the equation.
%
% When adding superscript or subscripts outside of brackets/braces, please
% group using {}.  For example, change "[U(D,E,\gamma)]^2" to
% "{[U(D,E,\gamma)]}^2". 
%
% Do not use \cal for caligraphic font.  Instead, use \mathcal{}
%
% % % % % % % % % % % % % % % % % % % % % % % % 
%
% Please contact latex@plos.org with any questions.
%
% % % % % % % % % % % % % % % % % % % % % % % %

\documentclass[10pt,letterpaper]{article}
\usepackage[top=0.85in,left=2.75in,footskip=0.75in]{geometry}

% amsmath and amssymb packages, useful for mathematical formulas and symbols
\usepackage{amsmath,amssymb}

% Use adjustwidth environment to exceed column width (see example table in text)
\usepackage{changepage}

% Use Unicode characters when possible
\usepackage[utf8x]{inputenc}

% textcomp package and marvosym package for additional characters
\usepackage{textcomp,marvosym}

% cite package, to clean up citations in the main text. Do not remove.
\usepackage{cite}

% Use nameref to cite supporting information files (see Supporting Information section for more info)
\usepackage{nameref,hyperref}

% line numbers
\usepackage[right]{lineno}

% ligatures disabled
\usepackage{microtype}
\DisableLigatures[f]{encoding = *, family = * }

% color can be used to apply background shading to table cells only
\usepackage[table]{xcolor}

% array package and thick rules for tables
\usepackage{array}

% package for vector symbols (e.g. \\vv)
\usepackage{esvect}

% new command for normalization (e.g. \norm{a} is |a|)
\newcommand{\norm}[1]{\left\lVert#1\right\rVert}

% new command for vectors
\renewcommand{\vector}[1]{\mathbf{#1}}

% new command for unit vectors
\newcommand{\uvector}[1]{\mathbf{\hat{#1}}}

% create "+" rule type for thick vertical lines
\newcolumntype{+}{!{\vrule width 2pt}}

% create \thickcline for thick horizontal lines of variable length
\newlength\savedwidth
\newcommand\thickcline[1]{%
  \noalign{\global\savedwidth\arrayrulewidth\global\arrayrulewidth 2pt}%
  \cline{#1}%
  \noalign{\vskip\arrayrulewidth}%
  \noalign{\global\arrayrulewidth\savedwidth}%
}

% \thickhline command for thick horizontal lines that span the table
\newcommand\thickhline{\noalign{\global\savedwidth\arrayrulewidth\global\arrayrulewidth 2pt}%
\hline
\noalign{\global\arrayrulewidth\savedwidth}}


% Remove comment for double spacing
%\usepackage{setspace} 
%\doublespacing

% Text layout
\raggedright
\setlength{\parindent}{0.5cm}
\textwidth 5.25in 
\textheight 8.75in

% Bold the 'Figure #' in the caption and separate it from the title/caption with a period
% Captions will be left justified
\usepackage[aboveskip=1pt,labelfont=bf,labelsep=period,justification=raggedright,singlelinecheck=off]{caption}
\renewcommand{\figurename}{Fig}

% Use the PLoS provided BiBTeX style
\bibliographystyle{plos2015}

% Remove brackets from numbering in List of References
\makeatletter
\renewcommand{\@biblabel}[1]{\quad#1.}
\makeatother

% Leave date blank
\date{}

% Header and Footer with logo
\usepackage{lastpage,fancyhdr,graphicx}
\usepackage{epstopdf}
\pagestyle{myheadings}
\pagestyle{fancy}
\fancyhf{}
\setlength{\headheight}{27.023pt}
\lhead{\includegraphics[width=2.0in]{PLOS-submission.eps}}
\rfoot{\thepage/\pageref{LastPage}}
\renewcommand{\footrule}{\hrule height 2pt \vspace{2mm}}
\fancyheadoffset[L]{2.25in}
\fancyfootoffset[L]{2.25in}
\lfoot{\sf PLOS}

%% Include all macros below

\newcommand{\lorem}{{\bf LOREM}}
\newcommand{\ipsum}{{\bf IPSUM}}

%% END MACROS SECTION


\begin{document}
\vspace*{0.2in}

% Title must be 250 characters or less.
\begin{flushleft}
{\Large
\textbf\newline{Calibration of a Murine Virtual Reality System}
}
\newline
% Insert author names, affiliations and corresponding author email (do not
% include titles, positions, or degrees).
\\
James W. Bridgewater\textsuperscript{1,2},
Reuben H. Fan\textsuperscript{1},
Dora E. Angelaki\textsuperscript{1},
Xaq S. Pitkow\textsuperscript{1,2*}
\\
\bigskip
\textbf{1} Department of Neuroscience, Baylor College of Medicine, Houston, TX,
USA
\\
\textbf{2} Department of Electrical and Computer Engineering, Rice University,
Houston, TX, USA
\\
\bigskip

% Insert additional author notes using the symbols described below. Insert
% symbol callouts after author names as necessary.
% 
% Remove or comment out the author notes below if they aren't used.
%
% Primary Equal Contribution Note
%\Yinyang These authors contributed equally to this work.

% Additional Equal Contribution Note
% Also use this double-dagger symbol for special authorship notes, such as
% senior authorship.
%\ddag These authors also contributed equally to this work.

% Current address notes
%\textcurrency Current Address: Dept/Program/Center, Institution Name, City,
%State, Country % change symbol to "\textcurrency a" if more than one current
%address note \textcurrency b Insert second current address \textcurrency c
%Insert third current address

% Deceased author note
%\dag Deceased

% Group/Consortium Author Note
%\textpilcrow Membership list can be found in the Acknowledgments section.

% Use the asterisk to denote corresponding authorship and provide email address
% in note below.
* xaq@rice.edu

\end{flushleft}
% Please keep the abstract below 300 words
\section*{Abstract}
For a visual virtual reality system to be effective, the projected visual scene
must engage the vast majority of the subject's visual field.  This is
challenging with rodents due to their exceptionally large visual coverage.
Several research groups have successfully developed virtual reality systems
capable of immersing rodents in a virtual visual world using a single projector
and a single non-planar mirror to project images to a curved screen. A major
challenge for implementing these systems is correcting the distortion of the
visual stimulus.  The ability to confidently evaluate and robustly minimize the
errors in the display of visual stimuli is key to successfully utilizing a
rodent virtual reality system. Here we detail a calibration technique that is
relatively simple to reproduce that minimizes visual display errors for a
virtual reality system that utilizes a partial sphere screen, a partial sphere
mirror, and a single projector.


% Please keep the Author Summary between 150 and 200 words
% Use first person. PLOS ONE authors please skip this step. 
% Author Summary not valid for PLOS ONE submissions.   
%\section*{Author Summary}

\linenumbers

% Use "Eq" instead of "Equation" for equation citations.
\section*{Introduction}
For decades, the study of rodent behavior has been performed using physical
arenas and mazes \cite{Moser2008RodentSpaceRev}.  Thus, the physical dimensions
of the laboratory space dictated the extent of rodent exploration. With the
rise of rodent virtual reality (VR) systems, the creation virtual spaces
unbound by the physical lab space became possible. In addition, rodents can be
sufficiently stabilized inside the VR systems to allow for stable neural
recording using electrophysiological or imaging methods
\cite{harvey2009intracellular, schmidt2013cellular, aronov2014engagement,
keller2012sensorimotor, ayaz2013locomotion, saleem2013integration}.  The
combination of these methodologies has revealed novel insights to how neurons
respond during active navigation \cite{harvey2009intracellular,
schmidt2013cellular, aronov2014engagement, keller2012sensorimotor,
ayaz2013locomotion, saleem2013integration}.

A major challenge in designing a virtual world for rodents is adequately
engaging their visual system. The visual field of mice and rats covers nearly
the entire half sphere above the horizon and a large portion of the half sphere
below it \cite{hughes1977topography, wagor1980retinotopic, schuett2002mapping,
wang2007area}. Current systems address this challenge with multiple visual
displays \cite{keller2012sensorimotor, ayaz2013locomotion,
saleem2013integration} or use non-planar mirrors to increase the visual field
covered by a single projector \cite{harvey2009intracellular,
schmidt2013cellular, aronov2014engagement}. The difficulty in using multiple
displays is the presence of bezels.  The use of multiple displays has, in some
cases, resulted in behavior which suggests that the animal is navigating the
real space occupied by the displays rather that the virtual space shown on them
\cite{holscher2005rats}. An alternative method is to employ non-planar mirrors
to increase the effective coverage of single projector
\cite{harvey2009intracellular, schmidt2013cellular, aronov2014engagement}.  The
most well-known examples use custom mirrors and projection surface shapes that,
while not prohibitive, adds a layer of complexity increasing the barrier of
entry \cite{holscher2005rats, harvey2009intracellular}. Alternatively, one can
use a simple spherical mirror and spherical dome, making VR construction more
accessible \cite {schmidt2013cellular}. Ultimately, the challenge of using
non-planar mirrors is correcting the distortion of the projected visual
stimuli. Correcting of the image distortion, in principle, requires an image
modification by solving a straight-forward 3-dimensional trigonometry problem
based on known specifications of the projector, mirror, projection surface, and
the observer’s location. While conceptually simple, in practice, correcting for
the visual distortions from non-planar mirrors can be challenging.

To achieve the best results, the system measurements are ideally known to
sub-millimeter accuracy. This is difficult to accomplish via measurement
because some elements of the system, like the projector's focal point, are not
directly observable. Thus, to overcome these issues and achieve a highly
accurate reproduction of viewing directions in our rodent VR system, we have
developed a technique that utilizes a 3D printed calibration device to project
spots of light in 37 known viewing directions relative the location of the
observer's head. By mapping projected points in the projector image to known
viewing directions displayed by the 3D printed device, it is possible to find
values of the virtual reality system's geometric parameters (see Methods for
list of parameters) that produce accurate calculated viewing directions by
minimizing the difference between the known and calculated viewing directions.
Once the system’s geometric parameters are known, an image can be modified such
that projection of the modified image in the virtual reality system produces
the original image. We will describe this method in detail. 

\section*{Materials and Methods}


\subsection*{System Geometry}

\subsection*{Viewing Direction Calculation}

Figure \ref{fig:SysScheme} is a depiction of our model VR system. Our objective
for the visual stimuli in this system is to cover ~ 270\textdegree{} horizontal
field of view, from 135\textdegree{} left of center to 135\textdegree{} right
of center, and a 90\textdegree{} vertical field of view, from 20\textdegree{}
below horizontal to 70\textdegree{} above horizontal. 

\begin{figure}[h]
\centering
\includegraphics[width=\textwidth]{SysScheme.pdf}
\caption{{\bf System schematic.} In our schematic, the configuration is set to
        emphasize the flexibility we have in the positioning of the point-A
        relative to the mirror, projector, and dome. Point-A is the observer's
        head location. Point-D is the origin of the sphere the projection dome
        is cut from. Point-E represents a location on the inner surface of the
        projection dome.  Point-M is defined as the center of the sphere the
        mirror is cut from and is used as the origin point for our entire
        system. Point-N is the location that light from focal point from the
        projector, defined as point-P, hits.}

\label{fig:SysScheme}
\end{figure}
Our model system uses a half sphere to display projected images and a quarter
sphere mirror to reflect an image coming from a single projector Fig.
\ref{fig:SysScheme}. The animal’s head position is represented as a small ball
sitting on a Styrofoam ball that is floating on a cushion of air similar to
other frictionless treadmill systems (REF). A three dimensional unit vector is
used to represent the viewing direction from the animal (Fig. 1 point-A) to a
point on the spherical screen (Fig. 1 point-E). The viewing direction
corresponding to a point in the projector image is calculated using the
following three steps:

\begin{enumerate}

\item
calculate the three dimensional Cartesian coordinates of point-N where light
from a given point in the projector image strikes the mirror,

\item
calculate the coordinates of point-E where the light reflected from the mirror
at point N strikes the screen,

\item
calculate the vector from the animal's position at point A to point E on the
screen and normalize it to get a unit vector.

\end{enumerate}

\newpage
\subsection*{Calculating the Coordinates of Point N}
In Fig.~1 Consider each side of the triangle consisting of points M, P, and N:

\begin{itemize}

\item
the vector from M to P can be calculated from the system's parameters,

\item
the length of the vector from M to N is equal to the radius of the mirror,
but its direction is unknown,

\item
the direction of the vector from P to N can be calculated from the system's
parameters for any given point in the projector image, but its length is
unknown.

\end{itemize}





Step 3 is to
use the triangle defined by point-A, point-E, and point-M to calculate the
vector from point-A to point-E and normalize it to get a unit vector that
represents the viewing direction corresponding to the given point in the
projector's image.  Note that in our method, we have chosen point-M to be the
origin of the system.

\begin{figure}[h]
\centering
\includegraphics[width=\textwidth]{FRUS.pdf}
\caption{{\bf Frustrum.}
The pixel position from the projected image are represented by the variables u and v. The resolution of our projector is 1280x720. The image spread (and height) and rise are related to the throw distance.}
\label{fig:FRUS}
\end{figure}
This calculation requires quantitative values for several geometric properties
of this virtual reality system. Specifically the location of the projector's
focal point, the aspect ratio of its image (Fig. \ref{fig:FRUS}), and an angle
quantifying how much the image spreads out horizontally as it gets farther from
the projector and a slope describing how far the image rises as it gets farther
from the projector (Fig. \ref{fig:FRUS}). Also required are the radius of the
spherical mirror, the radius of the hemispherical screen, the location of the
screen's center and the location of the animal's head relative to the projector
focal point, point-P (Fig. \ref{fig:SysScheme}). While many of these parameters
can be theoretically measured or provided, obtaining a sufficiently accurate
values for all of these parameters is challenging.  Particularly challenging is
that the focal point of the projector, the center of the hemispherical screen
and the center of the spherical mirror (the origin of our chosen coordinate
system) are not directly observable. We therefore developed the calibration
method described in the next section to find values of these parameters which
result in accurate calculated viewing directions.

\begin{eqnarray}
\label{eq:1}
  \vector{d}(r,c) = \frac{\vector{a}(r,c)}{\norm{\vector{a}(r,c)}}
\end{eqnarray}

\begin{eqnarray}
\label{eq:2}
  \vector{a}(r,c) = -\frac{w_l}{w_p}(c + 0.5 - 0.5w_p)\uvector{x} -\uvector{y} 
   + \frac{h_l}{h_p}(r + 0.5 - 0.5h_p)\uvector{z} + z_{off}\uvector{z}
\end{eqnarray}

\begin{eqnarray}
\label{eq:3}
  \vv{MN}(r,c) = \vv{MP} + \vv{PN}(r,c)
\end{eqnarray}

\begin{eqnarray}
\label{eq:4}
  \vv{PN}(r,c) = min(\frac{-2\vv{MP}\cdot\vector{d}(r,c)\pm
        \sqrt{(2\vv{MP}\cdot\vector{d}(r,c))^2 - 4(\norm{\vv{MP}}^2 -
        r_m^2)}}{2})\vector{d}(r,c)
\end{eqnarray}

\begin{eqnarray}
\label{eq:5}
  \vector{r}(r,c) = \frac{\vector{q}(r,c)}{\norm{\vector{q}(r,c)}}
\end{eqnarray}

\begin{eqnarray}
\label{eq:6}
  \vector{q}(r,c) = \vv{PN}(r,c) -
        2(\vv{PN}(r,c)\cdot\vector{n}(r,c))\ \vector{n}(r,c)
\end{eqnarray}

\begin{eqnarray}
\label{eq:7}
  \vector{n}(r,c) = \frac{\vv{MN}(r,c)}{\norm{\vv{MN}(r,c)}}
\end{eqnarray}

\begin{eqnarray}
\label{eq:8}
  \vv{NE}(r,c) = max(\frac{-2(\vv{MN} - \vv{MD})\cdot\vector{r}(r,c)\pm
        \sqrt{(2(\vv{MN} - \vv{MD})\cdot\vector{r}(r,c))^2 - 4(\norm{(\vv{MN} -
        \vv{MD})}^2 - r_d^2)}}{2})\vector{r}(r,c)
\end{eqnarray}

\begin{eqnarray}
\label{eq:9}
  \vector{v}(r,c) = \frac{\vv{AE}(r,c)}{\norm{\vv{AE}(r,c)}}
\end{eqnarray}

\begin{eqnarray}
\label{eq:10}
  \vv{AE}(r,c) = \vv{MN}(r,c) + \vv{NE}(r,c) - \vv{MA}
\end{eqnarray}

\begin{eqnarray}
\label{eq:lawofcosines}
  \norm{\vector{m}}^2 = \norm{\vector{i}}^2 + \norm{\vector{p}}^2 -
  2\norm{\vector{p}}\norm{\vector{i}}cos(\phi) 
\end{eqnarray}

\begin{eqnarray}
\label{eq:dotproduct}
  -\vector{p} \cdot \vector{i} = \norm{\vector{p}} \norm{\vector{i}} cos(\phi)
\end{eqnarray}

\begin{eqnarray}
\label{eq:dotproduct2}
  2\vector{p} \cdot \vector{i} = -2\norm{\vector{p}} \norm{\vector{i}} cos(\phi)
\end{eqnarray}

\begin{eqnarray}
\label{eq:lawofcosines2}
  \norm{\vector{m}}^2 = \norm{\vector{i}}^2 + \norm{\vector{p}}^2
  + 2\vector{p} \cdot \vector{i}
\end{eqnarray}

\begin{eqnarray}
\label{eq:lawofcosines3}
  0  = \norm{\vector{i}}^2 + 2\vector{p} \cdot \vector{i}  
  + \norm{\vector{p}}^2 -\norm{\vector{m}}^2
\end{eqnarray}

\begin{eqnarray}
\label{eq:quadratic}
  0  = \norm{\vector{i}}^2 + 2\vector{p} \cdot
  \frac{\vector{i}}{\norm{\vector{i}}} \norm{\vector{i}}
  + \norm{\vector{p}}^2 -\norm{\vector{m}}^2
\end{eqnarray}

\begin{eqnarray}
\label{eq:1}
  \vector{d}(r,c) = \frac{\vector{i}}{\norm{\vector{i}}}
\end{eqnarray}

\begin{eqnarray}
\label{eq:quadratic2}
  0  = \norm{\vector{i}}^2 + 2\vector{p} \cdot
  \vector{d}(r,c) \norm{\vector{i}}
  + \norm{\vector{p}}^2 -\norm{\vector{m}}^2
\end{eqnarray}

\begin{eqnarray}
\label{eq:quadraticsoln}
  \norm{\vector{i}} = \frac{-2\vector{p} \cdot \vector{d}(r,c) + 
  \sqrt{(2\vector{p} \cdot \vector{d}(r,c))^2
  - 4(\norm{\vector{p}}^2 - \norm{\vector{m}^2})}}{2}
\end{eqnarray}






(Description of equation above missing)

The aspect ratio of the image from the projector is provided by the
manufacturer. Here we used a DepthQ projector (1280x720 resolution).  The image
rise for our projector is estimated 15 percent of the image height at a given
throw distance. The image height is found by dividing the throw distance by the
factor 1.39 for the minimum magnification setting. The image width is computed
from the aspect ratio setting (16:9). The quarter sphere mirror is a standard
18-inch diameter safety mirror. The dome projection surface was custom made was
designed to have a radius of 70 cm. We leveled the projector and placed it at a
distance such that at the minimum magnification of the projector, the image
being displayed (before any correction) was qualitatively crisp. We then
proceeded to perform our custom calibration process.

\subsection*{Calibration Process}
One of the keys to making an accurate image from the perspective of the
observer is properly correcting for image distortion due to the use of a curved
mirror to reflect the projected image. We have developed a method for image
correction that is relatively simple to replicate. We use a 3D printed
calibration device (Fig. \ref{fig:3Ddev}) that projects 37 spots of light in
known directions, four rows of nine spots plus one spot straight overhead. Each
row has 9 spots located at yaw angles of -120, -90, -60, -30, 0, 30, 60, 90,
and 120 degrees where 0 degrees is straight ahead, negative angles refer to the
left visual hemifield and positive angles refer to locations in the right
visual hemifield. The four rows are at pitch angles of -15, 0, 30, and 60
degrees where negative angles referring to the lower visual field and positive
angles the upper visual field. This calibration device is placed at the
animal's head location inside the virtual reality system. The light source is a
diffuse light provided by a simple flashlight without any casing to confine the
light beams. The result is circles with a radius of approximately 1 visual
degree (Fig. \ref{fig:Dots}). To calibrate our image, we projected green dots
in grid-like fashion and allow those points fall onto the dome projection
surface. We then manually move each dot in a pixel wise fashion until each dot
falls into the proper corresponding spot projected by the calibration device.
We then save the pixel location of each dot emanating from the projector. This
file constitutes a mapping between 37 points in the projector image and 37
known viewing directions which can be used to find virtual reality system
geometry parameter values which produce accurate viewing directions. These
parameter values are found by minimizing the sum of the square L2 norms of the
vectors which constitute the differences between the calculated and known
viewing directions. The values we are particularly interested is are the y and
z coordinates of the dome, animal, and projector. The x coordinate is assumed
to be the same meaning that the center of the dome, animal, and projector image
center all fall on the same plane. This is done through careful construction. 

\begin{figure}[h]
\centering
\includegraphics[width=\textwidth]{3Dcaldev.pdf}
\caption{{\bf 3D printed calibration device.} A. Rendering of the 3D device. B.
        3D printed device without flashlight. C. 3D printed device with
        flashlight.}

\label{fig:3Ddev}
\end{figure}
\subsection*{Image Modification}

The result of the minimization routine provides us with the parameters
necessary to perform a pixel wise mapping from the projector to the appropriate
location on the dome.  The relationship between these two is given by:

INSERT EQUATION 2.

This relationship permits the user to create an image in a program such as
OpenGL and then warp the image such that once it is displayed on the dome after
reflecting off the mirror, the displayed image appears undistorted to the
observer. Figure 5 provides examples of corrected and uncorrected images and is
described in the Results section. 


%% For figure citations, please use "Fig" instead of "Figure".
%Nulla mi mi, Fig~\ref{fig1} venenatis sed ipsum varius, volutpat euismod diam.
%
%% Place figure captions after the first paragraph in which they are cited.
%\begin{figure}[!h]
%\caption{{\bf Bold the figure title.}
%Figure caption text here, please use this space for the figure panel
%descriptions instead of using subfigure commands. A: Lorem ipsum dolor sit
%amet. B: Consectetur adipiscing elit.}
%\label{fig1}
%\end{figure}

% Results and Discussion can be combined.
\section*{Results}

\subsection*{Finding the geometric values}

An important step for our method is obtaining quantitative values for several
geometric properties. As previously specified in the Materials and Methods, we
are particularly interested in mirror radius, the dome radius, and the y and z
position of the animal, dome, projector relative the mirror point-P (Fig. 1).
Our dome was designed to be approximately 70 cm in radius which is the value
our method provided. The mirror diameter is reported as 18 in, our algorithm
provided one of 17.7 in (1.67 percent smaller than expected). We estimated the
throw distance to be approximately 79.5 cm which would yield a rise position of
approximately 8.57 cm. The final rise position determined by our algorithm was
8.4 cm (2.1 percent smaller). While none of these values are hugely inaccurate,
the consistent difference does illustrate how difficult it is to find exact
values. While having hand measurements of these values is not necessary, they
are useful because they provide realistic constraints on our algorithm rather
than do a full exhaustive search. 

\begin{figure}[h]
\centering
\includegraphics[width=\textwidth]{Dots.pdf}
\caption{{\bf Calibration dots.}
The circles of light come from the 3D calibration device (see Fig. \ref{fig:3Ddev}). The green dots are the position after the Calibration process described in the Materials and Methods}
\label{fig:Dots}
\end{figure}
\subsection*{Image Modification}
Because we focus on only a few parameters of interest to simplify the
procedure, there are a few areas where the resulting projecting could be
improved. In Fig. \ref{fig:Dots}, the green dots are the final positions after
the calibration has been applied. Note that the green dots are not perfectly
centered. However the radius of each of the projected circles from the 3D
calibration device is approximately 1 visual degree. Thus as long as the green
dots fall within these circles they are very close to the ideal positions. The
reasons why they do not fall into the ideal positions are in part due to the
variables we do not allow to vary. We expand on this point in the Discussion.

\subsection*{Image display}

To illustrate the effectiveness of our method, we provide a set of corrected,
non-corrected image, and native image in Fig. \ref{fig:Compare}. The corrected
and non-corrected images were taken by a webcam inside the VR system from the
perspective of a viewer. GET A SET AND DESCRIBE. 

\begin{figure}[h]
\centering
\includegraphics[width=\textwidth]{original_projected.pdf}
\caption{{\bf Original and Projected images.} A/B. Original images. C/D. Image
        captures from a webcam placed at the location of the observer's head
        after projecting them onto the spherical projection surface.}

\label{fig:Compare}
\end{figure}

\subsection*{Image Modification}
Also notable is that we are able to perform these correction in a situation
where viewing distance of the observer is not the radius of the dome in all
viewing directions.  In our configuration, our observer is raised relative the
plane that marks the displayable surface of our half sphere. Our algorithm
places the animal 18.5 cm higher than this plane and produce good quality
images to the human eye. This has the same desired effect that pitching the
dome forward to provide some lower visual field coverage
\cite{schmidt2013cellular} or constructing a toroidal screen to encompass the
lower visual field \cite{holscher2005rats, harvey2009intracellular}. By raising
the animal, we are able to provide 20+ degrees of the lower visual field
without having to add any additional surface. 

% Place tables after the first paragraph in which they are cited.
%\begin{table}[!ht]
%\begin{adjustwidth}{-2.25in}{0in} % Comment out/remove adjustwidth environment if table fits in text column.
%\centering
%\caption{
%{\bf Table caption Nulla mi mi, venenatis sed ipsum varius, volutpat euismod diam.}}
%\begin{tabular}{|l+l|l|l|l|l|l|l|}
%\hline
%\multicolumn{4}{|l|}{\bf Heading1} & \multicolumn{4}{|l|}{\bf Heading2}\\ \thickhline
%$cell1 row1$ & cell2 row 1 & cell3 row 1 & cell4 row 1 & cell5 row 1 & cell6 row 1 & cell7 row 1 & cell8 row 1\\ \hline
%$cell1 row2$ & cell2 row 2 & cell3 row 2 & cell4 row 2 & cell5 row 2 & cell6 row 2 & cell7 row 2 & cell8 row 2\\ \hline
%$cell1 row3$ & cell2 row 3 & cell3 row 3 & cell4 row 3 & cell5 row 3 & cell6 row 3 & cell7 row 3 & cell8 row 3\\ \hline
%\end{tabular}
%\begin{flushleft} Table notes Phasellus venenatis, tortor nec vestibulum
%	mattis, massa tortor interdum felis, nec pellentesque metus tortor nec
%	nisl. Ut ornare mauris tellus, vel dapibus arcu suscipit sed.
%\end{flushleft}
%\label{table1}
%\end{adjustwidth}
%\end{table}


%PLOS does not support heading levels beyond the 3rd (no 4th level headings).
%\subsection*{\lorem\ and \ipsum\ Nunc blandit a tortor.}
%\subsubsection*{3rd Level Heading.} 
%
%\begin{enumerate}
%	\item{react}
%	\item{diffuse free particles}
%	\item{increment time by dt and go to 1}
%\end{enumerate}
%
%
%\subsection*{Sed ac quam id nisi malesuada congue.}
%
%
%\begin{itemize}
%	\item First bulleted item.
%	\item Second bulleted item.
%	\item Third bulleted item.
%\end{itemize}

\section*{Discussion}

Challenges
One of the challenges of our method, as presented, is that we assume the center
of the observer’s head, the quarter sphere mirror, and the center of the
projected image all fall onto a single plane. In practice, this is quite
difficult to achieve. We used a plumb line and displayed a center line and
manually moved the individual components until everything was bisected by the
plumb line. In theory we add this to the variable that we search for. But
increasing the number of variables we search for comes with its own concerns.
The inclusion of a third position parameters for the dome, the projector, and
the observer’s head increases the potential for local minima that have to be
avoided. While this is not likely to be too computationally expensive, the
addition of this parameter can make arriving at a values unreasonably
difficult. Thus we chose the difficulties in manually aligning these centers to
reduce the number of variables that need to be considered in our software. 

A second challenge is shadowing. In our configuration, the position of the
animal and platform used to support the animal will casts a shadow preventing
any images from being displayed. This is a design challenge that ideally falls
within the 90 degree wedge behind the rodent’s head which is close to full
horizontal visual field coverage (REF). This will be difficult depending on
what materials are used to support the chosen rodent. For rats the support
system necessary will cast a larger shadow than one for a mouse. However amount
of shadow cast can be mitigated, but not avoided, with clever design of the
supporting mechanism. Once the supporting mechanism is set, the amount of
shadowing can be assessed.

A final challenge is more of a consideration when designing a system similar to
that presented here. It is not advisable to attempt using a 9 in diameter
quarter sphere mirror with our dome. Similarly, it is not advisable to use an
18 in diameter quarter sphere mirror with a dome of 35 cm radius. In the
former, it would be difficult to find positions that would allow one to cover
enough of the dome surface to achieve adequate visual field coverage. The
latter would not have this issue but would present practical challenges given
the size of the mirror relative to the size of the projection dome.  

Benefits
The main benefit of our method is that it allows for the user to optimize the
visual scene for a large number of potential components and positions of the
animal.  While we have the space to house a large 70 cm radius half sphere.
This might not be the case for all. However our method allows others to choose
sizes that better fit the constraints of their lab and provides them with a
method to correct for visual distortions. In addition, we could theoretically
use less half of sphere and our method would still work. This flexibility of
design is permissible because we allow for the observer to be in any position
inside the dome. In our example we move the position of the observer’s head
18.5 cm above point-D (Fig. 1). This allows us to provide 20+ degrees of the
lower visual field without needing to add additional projection space. Past
examples have placed the viewer at point-D (see REF). The reason for this is
similar to our aligning the observer’s head, the mirrors center, the dome’s
center, and the projected image center. By centering the mouse and the
deflecting mirror, several parameters are simplified. One of the parameters is
the viewing distance to all points on the screen are defined by the radius of
the dome. This would reduce the complexity of the distortion correction
considerably. However the disadvantage of this is such a configuration
necessarily sacrifices visual field coverage unless the dome is constructed in
a fashion that extends the projection surface. Another tactic to gaining more
visual field is to use cones or toroidal screens with custom mirrors
\cite{holscher2005rats, harvey2009intracellular}, to use a cone shape
\cite{aronov2014engagement}, or simple boxes \cite{aghajan2015impaired}.
Regardless of the specifics components used to construct an immersive VR
system, in the end, the problem still remains the same, how to calibrate the
visual scene. By allowing for different geometric configurations providing
parts of the lower visual field of locations behind the head are a simple
matter of simply moving the observer into the dome a certain distance. Because
we use simple components that are more widely accessible, our method should be
more accessible by a wider array of experimenters. However should other shapes
be used, our method is theoretically adjustable to fit a wider array of
potential system designs.

\section*{Conclusion}

information, see \nameref{S1_Appendix}.

\section*{Supporting Information}

% Include only the SI item label in the paragraph heading. Use the
% \nameref{label} command to cite SI items in the text.

\paragraph*{S1 Fig.}
\label{S1_Fig}
{\bf Bold the title sentence.} Add descriptive text after the title of the item
(optional).

\paragraph*{S2 Fig.}
\label{S2_Fig}
{\bf Lorem Ipsum.} Maecenas convallis mauris sit amet sem ultrices gravida.
Etiam eget sapien nibh. Sed ac ipsum eget enim egestas ullamcorper nec euismod
ligula. Curabitur fringilla pulvinar lectus consectetur pellentesque.

\paragraph*{S1 File.}
\label{S1_File}
{\bf Lorem Ipsum.}  Maecenas convallis mauris sit amet sem ultrices gravida.
Etiam eget sapien nibh. Sed ac ipsum eget enim egestas ullamcorper nec euismod
ligula. Curabitur fringilla pulvinar lectus consectetur pellentesque.

\paragraph*{S1 Video.}
\label{S1_Video}
{\bf Lorem Ipsum.}  Maecenas convallis mauris sit amet sem ultrices gravida.
Etiam eget sapien nibh. Sed ac ipsum eget enim egestas ullamcorper nec euismod
ligula. Curabitur fringilla pulvinar lectus consectetur pellentesque.

\paragraph*{S1 Appendix.}
\label{S1_Appendix}
{\bf Lorem Ipsum.} Maecenas convallis mauris sit amet sem ultrices gravida.
Etiam eget sapien nibh. Sed ac ipsum eget enim egestas ullamcorper nec euismod
ligula. Curabitur fringilla pulvinar lectus consectetur pellentesque.

\paragraph*{S1 Table.}
\label{S1_Table}
{\bf Lorem Ipsum.} Maecenas convallis mauris sit amet sem ultrices gravida.
Etiam eget sapien nibh. Sed ac ipsum eget enim egestas ullamcorper nec euismod
ligula. Curabitur fringilla pulvinar lectus consectetur pellentesque.

\section*{Acknowledgments}

This work was supported by NSF IOS-1450923, the Simons Foundation, and the
McNair Foundation.

\nolinenumbers

% Either type in your references using
% \begin{thebibliography}{}
% \bibitem{}
% Text
% \end{thebibliography}
%
% or
%
% Compile your BiBTeX database using our plos2015.bst
% style file and paste the contents of your .bbl file
% here.
% 

\bibliography{bibliography}


\end{document}

